
\section{Method}

\subsection{Overview}

SALT replaces each linear projection $W$ with a three-component parameterization:

\begin{lstlisting}
    Weight W $\in$ R^{d_out $$\times$ d_in}
            |
             
    +-------------------------+
    |  Truncated SVD          |
    |  W $\approx$ U_r $\Sigma$_r V_r^T      |
    +-----------+-------------+
                |
    +-----------+------------+
                              
    +--------------+   +-----------------+
    |  Buffers     |   |  Trainable      |
    |  U, S, Vt    |   |  $\gamma$, A, B        |
    |  (frozen)    |   |                 |
    +--------------+   +-----------------+
                |                 |
                +--------+--------+
                          
    +--------------------------------------+
    |  Forward:                             |
    |  W_eff = U diag(S   $\gamma$) V^T + A B     |
    |  y = W_eff x + b                      |
    +--------------------------------------+
\end{lstlisting}

\subsection{The SVD Skeleton}

Given a pretrained weight $W \in \mathbb{R}^{d_\text{out} \times d_\text{in}}$, 
compute its SVD:

\[
W = U \Sigma V^\top, \quad U \in \mathbb{R}^{d_\text{out} \times d_\text{out}}, 
\Sigma \in \mathbb{R}^{d_\text{out}}, V \in \mathbb{R}^{d_\text{in} \times d_\text{in}}
\tag{4.1}
\]

Keep the top-$r$:

\[
U_r = U_{:, :r}, \quad S_r = \Sigma_{:r}, \quad V_r = V_{:r, :}^\top
\tag{4.2}
\]

We store $U_r, S_r, V_r$ as \textbf{buffers} (non-trainable), typically in the 
model's dtype (FP16 or FP32).

\subsection{Singular Value Modulation}

We add a trainable scale $\gamma \in \mathbb{R}^r$, initialized to 1:

\[
\Sigma_\gamma = S_r \odot \gamma
\tag{4.3}
\]

This allows the model to amplify or suppress retained singular directions.

\textbf{Why is this useful?}
\begin{itemize}
\item The pretrained SVD is optimal in Frobenius norm but not for the downstream
  task.
\item $\gamma$ provides a small number of trainable parameters ($r$) that adjust
  the spectral profile.

\end{itemize}
\subsection{Low-Rank Residual}

We add a trainable low-rank residual $AB$:

\[
A \in \mathbb{R}^{d_\text{out} \times r}, \quad 
B \in \mathbb{R}^{r \times d_\text{in}}
\tag{4.4}
\]

\textbf{Initialization:}
\begin{itemize}
\item $A \sim \mathcal{N}(0, 0.01^2)$ (small random)
\item $B = 0$

\end{itemize}
So $AB = 0$ at initialization.

\textbf{Why a residual?} The SVD skeleton captures the top-$r$ directions. If 
these are insufficient, the residual can add capacity -- but starting from 
zero avoids disrupting the pretrained state.

\subsection{Combined Weight}

The effective weight is:

\[
W_\text{SALT} = U_r \cdot \text{diag}(S_r \odot \gamma) \cdot V_r + A B
\tag{4.5}
\]

\textbf{At initialization:} $W_\text{SALT} = U_r \text{diag}(S_r) V_r$ -- the 
rank-$r$ truncated SVD of $W_\text{pretrained}$.

\subsection{Forward Pass}

\subsubsection{Naive (Our Implementation)}

\begin{lstlisting}
def forward(self, x):
    adapted_S = self.S * self.gamma
    W_svd = (self.U * adapted_S.unsqueeze(0)) @ self.Vt  # reconstruct
    W_lora = self.A @ self.B                             # reconstruct
    W = (W_svd + W_lora).to(dtype=x.dtype)
    return F.linear(x, W, self.bias)                     # full matmul
\end{lstlisting}

\textbf{Cost.} $\mathcal{O}(d_\text{out} \cdot r \cdot d_\text{in} + d_\text{out} \cdot d_\text{in})$ 
per forward -- \textbf{worse than baseline} due to reconstruction.

\subsubsection{Fused (Future Work)}

The ideal forward avoids reconstructing W:

\begin{lstlisting}
def forward(self, x):
    # 1. Project to rank space
    z = F.linear(x, self.Vt)                        # (B, L, r)
    # 2. Scale by S * gamma
    z = z * (self.S * self.gamma)
    # 3. Project back
    y_svd = F.linear(z, self.U)                     # (B, L, d_out)
    # 4. LoRA branch
    y_lora = F.linear(F.linear(x, self.B), self.A)
    return y_svd + y_lora + self.bias
\end{lstlisting}

\textbf{Cost.} $\mathcal{O}(r(d_\text{out} + d_\text{in}))$ -- same as low-rank W.

\textbf{We do NOT implement the fused version.} The paper documents the naive 
version and its limitation.

\subsection{Parameter Counting}

\begin{table}[H]
\centering
\small
\begin{tabular}{llll}
\toprule
\textbf{Component} & \textbf{Shape} & \textbf{Trainable?} & \textbf{Params} \\
\midrule
$U_r$ & $(d_\text{out}, r)$ \& No (buffer) & $d_\text{out} \cdot r$ \\
$S_r$ & $(r,)$ \& No (buffer) & $r$ \\
$V_r$ & $(r, d_\text{in})$ \& No (buffer) & $r \cdot d_\text{in}$ \\
$\gamma$ & $(r,)$ \& Yes & $r$ \\
$A$ & $(d_\text{out}, r)$ \& Yes & $d_\text{out} \cdot r$ \\
$B$ & $(r, d_\text{in})$ \& Yes & $r \cdot d_\text{in}$ \\
bias & $(d_\text{out},)$ \& Yes & $d_\text{out}$ \\
\textbf{Total} &  &  & $2r(d_\text{out} + d_\text{in}) + 2r + d_\text{out}$ \\
\textbf{Full W} &  &  & $d_\text{out} \cdot d_\text{in}$ \\
\textbf{Ratio (768, r=16)} &  &  & \textbf{11.8$$\times$} \\
\bottomrule
\end{tabular}
\end{table}

\subsection{Weight Initialization}

\subsubsection{SVD Buffers}

For each projection, compute SVD of the pretrained weight and keep the 
top-$r$.

\subsubsection{Trainable Parameters}

\begin{itemize}
\item $\gamma \leftarrow \mathbf{1}$
\item $A \sim \mathcal{N}(0, 0.01^2)$
\item $B \leftarrow \mathbf{0}$
\item bias $\leftarrow$ pretrained bias

\end{itemize}
\subsection{Algorithm}

\textbf{Algorithm 1:} SALT initialization
\begin{lstlisting}
Input:  W $\in$ R^{d_out $$\times$ d_in}, bias b, rank r
Output: SALTLinear module

1: U, S, Vt $\leftarrow$ svd(W, full_matrices=False)
2: U_r $\leftarrow$ U[:, :r];  S_r $\leftarrow$ S[:r];  Vt_r $\leftarrow$ Vt[:r, :]
3: register_buffer("U", U_r)
4: register_buffer("S", S_r)
5: register_buffer("Vt", Vt_r)
6: $\gamma$ $\leftarrow$ ones(r)
7: A $\leftarrow$ randn(d_out, r) * 0.01
8: B $\leftarrow$ zeros(r, d_in)
9: bias $\leftarrow$ b
10: return SALTLinear(U_r, S_r, Vt_r, $\gamma$, A, B, bias)
\end{lstlisting}

\textbf{Algorithm 2:} SALT forward (naive)
\begin{lstlisting}
Input:  x $\in$ R^{B $$\times$ L $$\times$ d_in}
Output: y $\in$ R^{B $$\times$ L $$\times$ d_out}

1: S_$\gamma$ $\leftarrow$ S * $\gamma$
2: W_svd $\leftarrow$ (U * S_$\gamma$) @ Vt            (d_out, d_in)
3: W_lora $\leftarrow$ A @ B                     (d_out, d_in)
4: W $\leftarrow$ W_svd + W_lora
5: y $\leftarrow$ x @ W^T + bias
6: return y
\end{lstlisting}

\subsection{Complexity Analysis}

\begin{table}[H]
\centering
\small
\begin{tabular}{lll}
\toprule
\textbf{Operation} & \textbf{Time} & \textbf{Memory} \\
\midrule
Naive forward & $\mathcal{O}(BL(d_\text{out} d_\text{in} + r(d_\text{out} + d_\text{in})))$ & $\mathcal{O}(d_\text{out} d_\text{in})$ \\
Fused forward & $\mathcal{O}(BL r (d_\text{out} + d_\text{in}))$ & $\mathcal{O}(r(d_\text{out} + d_\text{in}))$ \\
Baseline & $\mathcal{O}(BL d_\text{out} d_\text{in})$ & $\mathcal{O}(d_\text{out} d_\text{in})$ \\
\bottomrule
\end{tabular}
\end{table}

\textbf{The naive forward has extra work} from reconstruction. The fused forward 
realizes the low-rank speedup but is not implemented in this work.

\subsection{Design Choices}

\subsubsection{Why SVD Buffers Instead of Trainable?}

Buffers avoid accidental modification. If $U, S, V$ were trainable, the model 
could drift arbitrarily far from the pretrained state. Freezing them ensures 
that changes are localized to $\gamma$ and $AB$.

\subsubsection{Why $\gamma$?}

A scalar per singular value provides a minimal but expressive modulation. 
It has only $r$ parameters (16 for r=16) -- negligible overhead.

\subsubsection{Why A ~ Small Random, B = 0?}

Standard LoRA init. Ensures $AB = 0$ at start, so the model matches the 
truncated SVD initialization.

\subsubsection{Why the Same Rank for SVD and Residual?}

Using a single rank parameter simplifies the API. In principle, SVD rank and 
residual rank could differ.
\begin{lstlisting}

---
\end{lstlisting}